\hypertarget{appendix-a6f2c4fe1996}{}
\section*{Eye movements and display-time gaze}
\section{Eye movements and display-time gaze}\label{eye-movements-and-display-time-gaze}

\textbf{Curriculum depth:} Introduced · \textbf{First seen:} XR01 · \textbf{Latest development:} XR01

\subsection{Definition and mechanism}\label{definition-and-mechanism}

Fixation, saccades, pursuit and VOR have different dynamics. Sampling rate differs from delivery/end-to-end delay. A rendering allocation should target visible display time, with phase-specific errors and uncertainty rather than a universal constant-velocity assumption.

\subsection{Prerequisites}\label{prerequisites}

\begin{itemize}
\tightlist
\item
  \hyperref[appendix-7d356c396d1a]{Visual angle and calibrated rays} --- Angular coordinates make sampling density and gaze-relative location physically meaningful.
\end{itemize}

\subsection{Used by}\label{used-by}

\begin{itemize}
\tightlist
\item
  \hyperref[appendix-70895739febb]{Spatiotemporal sampling and retinal motion}
\item
  \hyperref[appendix-5878b813b732]{Foveated rendering and uncertainty}
\end{itemize}

\subsection{Recall and next depth}\label{recall-and-next-depth}

Explain a gaze-landing failure despite low average angular error. Pending: tracking, fusion, prediction and timing pipelines. Written coverage is not a claim of learner mastery or accepted numbered-curriculum progress.

\subsection{Source and navigation}\label{source-and-navigation}

\href{Continuous\%20Notes/University\%20-\%20Worlds,\%20Images\%20and\%20Inference\%20-\%20Graphics,\%20Vision\%20and\%20Photography\#A\%20gaze\%20sample\%20needs\%20a\%20time\%20as\%20well\%20as\%20a\%20direction}{Expanded teaching explanation} · \href{}{Independent coverage (01-Foundations-I, appendix-fd090d982a15 § XR01 - Independent seeing and representation source)} · \href{}{Dependency Map (02-Foundations-II, appendix-4ad47bdd269a)}
